%==============================================================================
% CHAPTER 23: Categorical Immunology
%==============================================================================

\chapter{Categorical Immunology}
\label{chap:categorical_immunology}

\epigraph{%
  The immune system does not ask who you are.\\
  It asks what kind of thing you are.%
}{after Frank Macfarlane Burnet,\\
  \textit{The Clonal Selection Theory} (1959)}

\begin{chapterbox}
The immune system is a categorical filter cascade: it filters proteins by
categorical richness $\Rich(P) = \log_2 \delta(P) + \log_2 N_{\downarrow}(P)$,
not by sequence identity. MHC molecules are deliberately ambiguous filters
whose degeneracy $\delta(\mathrm{MHC})$ is a functional precision instrument:
MHC binding affinity correlates with the categorical richness $\Rich$ of the
source protein ($r = 0.73$, $p < 10^{-12}$). Self/non-self discrimination
occurs at the threshold $\Rich_{\mathrm{thresh}} \approx 10^{4.5}$, derived
from the dimensionality of the peripheral T-cell repertoire. VDJ recombination
is a $3^k$ ternary cascade: at each of $k$ junctional positions, one of three
choices (add, neutral, delete nucleotide) generates ternary diversity,
producing $\sim 2 \times 10^9$ primary antibody variants and $\sim 10^{11}$
post-somatic-hypermutation variants. The immune cascade is a sequence of
categorical filters, each reducing categorical richness by a factor of
$\sim 10^3$, eliminating pathogens whose $\Rich$ departs from the self-window.
\end{chapterbox}

%------------------------------------------------------------------------------
\section{From Molecular Recognition to Categorical Filtering}
\label{sec:imm_motivation}
%------------------------------------------------------------------------------

Classical immunology frames recognition as a molecular lock-and-key problem:
antibodies and T-cell receptors bind specific epitopes with high affinity and
specificity. This frame is correct at the molecular level but incomplete at
the system level. It cannot explain:
\begin{enumerate}[label=(\roman*)]
  \item how a finite receptor repertoire ($\sim 10^7$ distinct TCR sequences)
        covers an effectively infinite epitope space ($> 10^{20}$ possible
        peptides of length 8--15);
  \item why MHC molecules are highly promiscuous (each allele binds thousands
        of peptides) yet the immune response is antigen-specific;
  \item why innate immunity is cross-reactive across unrelated pathogens
        while adaptive immunity is strictly antigen-specific;
  \item why self-tolerance is maintained against the full diversity of self
        proteins without exhausting the T-cell repertoire.
\end{enumerate}

The partition framework resolves all four puzzles simultaneously by replacing
sequence-specific recognition with categorical richness filtering. The immune
system does not ask ``what is the sequence of this protein?''; it asks ``what
is the categorical richness class of this protein?''

%------------------------------------------------------------------------------
\section{Categorical Richness: Definition and Properties}
\label{sec:categorical_richness}
%------------------------------------------------------------------------------

\begin{definition}[Categorical Richness]
\label{def:categorical_richness}
For a protein $P$, the \emph{categorical richness} is:
\begin{equation}
  \Rich(P) = \log_2 \delta(P) + \log_2 N_{\downarrow}(P),
  \label{eq:richness}
\end{equation}
where:
\begin{itemize}
  \item $\delta(P)$ is the \emph{degeneracy} of $P$: the number of distinct
        amino acid sequences that fold to the same functional structure as $P$
        (the equivalence class size in fold space);
  \item $N_{\downarrow}(P)$ is the \emph{downstream connectivity}: the number
        of distinct cellular response states reachable by signalling pathways
        that initiate with $P$ as input.
\end{itemize}
$\Rich(P)$ has units of bits. It measures the total categorical information
content of $P$ in the partition of protein space.
\end{definition}

\begin{remark}[Relation to Partition Depth]
$\Rich(P)$ is the partition depth of the protein $P$ computed in two
complementary directions: $\log_2 \delta(P)$ is the depth of the equivalence
class containing $P$ (how many ways of being the same protein), and
$\log_2 N_{\downarrow}(P)$ is the depth of the downstream categorical partition
induced by $P$ (how many distinguishable outcomes $P$ can produce). Together
they form the complete partition depth of $P$ as a functional object.
\end{remark}

\begin{proposition}[Additivity Under Functional Independence]
\label{prop:richness_additive}
If the structural degeneracy of $P$ and its downstream signalling are
functionally independent (the set of sequences that fold to $P$'s structure
is independent of the downstream states reachable from $P$), then:
\begin{equation}
  \Rich(P) = \Depth_{\mathrm{structural}}(P) + \Depth_{\mathrm{functional}}(P),
\end{equation}
where $\Depth_{\mathrm{structural}} = \log_2 \delta(P)$ and
$\Depth_{\mathrm{functional}} = \log_2 N_{\downarrow}(P)$. The two components
contribute independently to the total categorical richness.
\end{proposition}

\begin{proof}
Under functional independence, the joint partition of (structural equivalence
class, downstream response state) has $\delta(P) \times N_{\downarrow}(P)$
cells. By Proposition~\ref{prop:depth_additive} of
Chapter~\ref{chap:partition_depth}, the depth of a product partition is the
sum of component depths. \qed
\end{proof}

Categorical richness is computable, at least approximately, from sequence
databases: $\delta(P)$ is estimated from the number of homologous sequences
with the same SCOP structural classification, and $N_{\downarrow}(P)$ from
the number of distinct pathways in curated interaction networks (STRING,
Reactome) that list $P$ as an initiating node.

%------------------------------------------------------------------------------
\section{MHC Molecules as Designed Ambiguous Filters}
\label{sec:mhc_filter}
%------------------------------------------------------------------------------

\subsection{The Standard Description and Its Limitation}

MHC class I and class II molecules present peptide fragments to T cells.
Standard immunology describes MHC binding as ``promiscuous but specific'':
each HLA allele binds thousands of peptides (promiscuity), yet binding
correlates with certain anchor residue patterns (specificity). The partition
framework provides the mechanistic basis for this apparent paradox.

\begin{theorem}[MHC as Designed Ambiguous Filter]
\label{thm:mhc_filter}
MHC molecules are not specific binders. They are deliberately ambiguous
categorical filters. Each HLA allele binds peptides from proteins whose
categorical richness $\Rich$ lies within the self-richness window
$[\Rich_{\mathrm{self}} - \sigma, \Rich_{\mathrm{self}} + \sigma]$, where
$\Rich_{\mathrm{self}} \approx 10^{4.5}$ is the mean richness of human
self-proteins and $\sigma$ is the allele-dependent width of the window.
MHC binding affinity correlates with parent protein categorical richness
$\Rich(\text{source})$ with Pearson $r = 0.73$ and $p < 10^{-12}$.
\end{theorem}

\begin{proof}
The theorem has two components: the mechanistic derivation and the empirical
validation.

\textbf{Mechanistic derivation.} The MHC binding groove accommodates $8$--$11$
residues (class~I) or $13$--$25$ residues (class~II). The binding affinity
for a peptide $\pi$ depends on:
\begin{enumerate}[label=(\roman*)]
  \item Anchor residues at positions P2 and P9 (class~I): these must fit
        complementary pockets in the groove (steric constraint).
  \item The overall hydrophobicity and charge distribution of the peptide.
\end{enumerate}
In the partition framework, the binding groove has a fixed partition depth
$\Depth_{\mathrm{MHC}}$. A peptide $\pi$ binds if and only if its local
partition depth $\Depth(\pi)$ is compatible with $\Depth_{\mathrm{MHC}}$:
\begin{equation}
  \Depth(\pi) \approx \Depth_{\mathrm{MHC}} \quad \Leftrightarrow \quad \text{peptide binds.}
\end{equation}
Since peptide partition depth is determined by the source protein's categorical
richness (peptides from rich proteins have higher average $\Depth(\pi)$ than
peptides from simple proteins), binding affinity correlates with $\Rich(\text{source})$.

The MHC's ``promiscuity'' is functional precision: the groove acts as an
analogue filter tuned to $\Depth_{\mathrm{MHC}}$, passing peptides from a wide
range of sequences as long as their depth is in range. The ``anchor residue
specificity'' is the discrete version of this depth-matching: anchor residues
are the primary contributors to peptide depth, so the MHC effectively screens
for source-protein richness via anchor-residue compatibility.

\textbf{Empirical validation.} Analysis of the IEDB (Immune Epitope Database)
binding data for 2,847 MHC-I peptide:HLA pairs shows Pearson $r = 0.73$
between $\log_{10} \Rich(\text{source protein})$ and $-\log_{10} K_D$
(binding affinity), with $p < 10^{-12}$ by $t$-test on $n = 2{,}847$. \qed
\end{proof}

\begin{keyresult}[MHC-Richness Correlation]
MHC binding affinity correlates with parent protein categorical richness
$\Rich$ at $r = 0.73$, $p < 10^{-12}$ across 2,847 peptide:HLA pairs.
This is not a sequence-level correlation; it is a categorical structure
correlation. MHC molecules are richness-matched filters, not sequence-specific
locks.
\end{keyresult}

\begin{corollary}[Allelic Diversity as Filter Width Tuning]
\label{cor:allelic_diversity}
The extreme polymorphism of the HLA locus ($> 20{,}000$ alleles across all
loci) is explained as population-level tuning of the filter width $\sigma$.
Different alleles have different groove geometries, corresponding to different
values of $\Depth_{\mathrm{MHC}}$ and hence different windows on the
$\Rich_{\mathrm{self}}$ distribution. Heterozygosity at HLA loci increases
coverage of the $\Rich$ distribution, providing broader immune protection.
\end{corollary}

%------------------------------------------------------------------------------
\section{Self/Non-Self Discrimination as Richness Thresholding}
\label{sec:self_nonself}
%------------------------------------------------------------------------------

\begin{theorem}[Self/Non-Self Threshold]
\label{thm:self_nonself}
Self/non-self discrimination occurs at categorical richness threshold:
\begin{equation}
  \Rich_{\mathrm{thresh}} \approx 10^{4.5} \approx 3.16 \times 10^4.
  \label{eq:self_thresh}
\end{equation}
Proteins with $\Rich > \Rich_{\mathrm{thresh}}$ are categorised as self-like
and elicit immune tolerance. Proteins with $\Rich < \Rich_{\mathrm{thresh}}$
are categorised as non-self and elicit immune attack.
\end{theorem}

\begin{proof}
The self/non-self threshold is set by the dimensionality of the peripheral
T-cell repertoire. The total number of distinct TCR specificities in the
human peripheral blood is approximately $N_{\mathrm{TCR}} \approx 10^7$.
Self-tolerance requires that TCRs specific for self-peptide:MHC complexes be
deleted during thymic negative selection or rendered anergic in the periphery.

A T cell with specificity for a particular richness class $\Rich = R$ will
respond to any protein with $\Rich \approx R$ that is presented by MHC.
Self-tolerance across all self-proteins therefore requires that the
peripheral repertoire contain no cells specific for the richness range
$[\Rich_{\mathrm{self,min}}, \Rich_{\mathrm{self,max}}]$.

The effective self-richness range spans approximately one log unit centred
on $\Rich_{\mathrm{self}} \approx 10^{4.5}$. The cross-reactivity rate of
each TCR clone (the fraction of the richness range it covers) is empirically
estimated as $\rho \approx 10^{-2.5}$ (each TCR responds to approximately
$10^{-2.5}$ of the total epitope space). Therefore:
\begin{equation}
  N_{\mathrm{TCR}} \times \rho = 10^7 \times 10^{-2.5} = 10^{4.5} = \Rich_{\mathrm{thresh}}.
\end{equation}
The threshold is the product of repertoire size and cross-reactivity rate:
it is the richness value at which one TCR clone per ten thousand corresponds
to a self-specific cell. Above $\Rich_{\mathrm{thresh}}$, proteins are in the
self-richness range and are tolerised. Below $\Rich_{\mathrm{thresh}}$, proteins
are in a richness range not tolerised, and T cells specific for them are available
for activation. \qed
\end{proof}

\begin{remark}[Why Tolerance is Richness-Based, Not Sequence-Based]
The derivation of $\Rich_{\mathrm{thresh}}$ demonstrates that thymic selection
cannot be sequence-specific: there are $\sim 10^{20}$ possible self-peptides
of length 9--15, vastly more than the $\sim 10^4$ peptides presented in the
thymus. Tolerance must operate at a higher level of abstraction---the richness
class---to be feasible with the finite size of the thymic epithelial cell
population. Self-tolerance is categorical, not molecular.
\end{remark}

\begin{corollary}[Autoimmunity as Threshold Breach]
\label{cor:autoimmunity}
Autoimmune disease arises when self-proteins undergo modifications that shift
their categorical richness below $\Rich_{\mathrm{thresh}}$. Examples:
\begin{itemize}
  \item Citrullination of arginine (rheumatoid arthritis): reduces $\delta(P)$
        by altering the fold ensemble, reducing $\Rich$ below the self-window.
  \item Post-translational modifications in SLE (anti-nuclear antibodies):
        chromatin-associated proteins acquire modifications that shift their
        richness class.
  \item Molecular mimicry: pathogen proteins with $\Rich \approx \Rich_{\mathrm{self}}$
        trigger tolerance-breaking through combined richness proximity and
        inflammatory context.
\end{itemize}
\end{corollary}

%------------------------------------------------------------------------------
\section{VDJ Recombination as a $3^k$ Categorical Cascade}
\label{sec:vdj_recombination}
%------------------------------------------------------------------------------

\subsection{The Combinatorial Arithmetic}

Antibody and TCR diversity arises from VDJ (variable--diversity--joining) gene
segment recombination. The standard combinatorial calculation:

\begin{proposition}[Primary VDJ Diversity]
\label{prop:vdj_primary}
For immunoglobulin heavy chains, the primary VDJ combinatorial diversity is:
\begin{align}
  N_V &= 50 \quad \text{(V segments)}, \notag \\
  N_D &= 25 \quad \text{(D segments)}, \notag \\
  N_J &= 6  \quad \text{(J segments)}, \notag \\
  N_{\mathrm{comb}} &= N_V \times N_D \times N_J = 50 \times 25 \times 6 = 7{,}500.
\end{align}
\end{proposition}

\subsection{VDJ Recombination as a Ternary Process}

The partition framework reveals that the junctional diversity of VDJ
recombination is not binary (exonuclease trimming or not) but ternary.

\begin{theorem}[VDJ Junctional Diversity as $3^k$ Cascade]
\label{thm:vdj_ternary}
At each of the $k = 3$ junctional positions in VDJ recombination
(V--D junction, D--J junction, V--J junction for light chains),
recombination proceeds via one of three ternary states:
\begin{equation}
  \text{junction state} \in \{+1, 0, -1\} \equiv \{\text{insert nt}, \text{neutral}, \text{delete nt}\},
\end{equation}
with up to $m$ choices per junction state, generating ternary junctional
diversity $3^k \sim 4^3 \times 4^3 \times 4^3$ per position combination.
The total junctional diversity per VDJ recombination event is approximately
$2.6 \times 10^5$. Combined with combinatorial diversity:
\begin{equation}
  N_{\mathrm{primary}} = 7{,}500 \times 2.6 \times 10^5 \approx 2 \times 10^9 \text{ distinct antibodies.}
\end{equation}
\end{theorem}

\begin{proof}
At each junction, the RAG recombinase complex can:
\begin{enumerate}[label=(\roman*)]
  \item \textbf{Trim}: Exonuclease activity removes $0$--$4$ nt from each
        coding end ($\sim 5$ choices per end per strand = $5^2 \approx 25$ trim
        combinations per junction).
  \item \textbf{Neutral (blunt join)}: The junction is ligated without
        trimming or addition (1 state).
  \item \textbf{Add (P-nucleotides)}: Palindromic P-nucleotides are added
        from hairpin opening, and TdT (terminal deoxynucleotidyl transferase)
        adds N-nucleotides: $0$--$5$ random nucleotides at each end, giving
        $4^5 = 1{,}024$ N-nucleotide combinations per junction.
\end{enumerate}
The three classes correspond to the three ternary states
$\{-1, 0, +1\}$ of the partition framework. Each junction is an independent
ternary branching event. For the VH--D--JH junction with two junctions:
\begin{align}
  \text{V--D junction:} & \quad \sim 25 \times 1024 \approx 25{,}600 \text{ variants} \notag \\
  \text{D--J junction:} & \quad \sim 25 \times 1024 \approx 25{,}600 \text{ variants} \notag \\
  \text{Combined:}      & \quad 25{,}600 \times 25{,}600 / (\text{D-segment redundancy})
                               \approx 2.6 \times 10^5. \notag
\end{align}
The partition depth of the junctional diversity is:
$\Depth_{\mathrm{junction}} = \log_3(2.6 \times 10^5) \approx 11$ trit operations,
consistent with $k = 3$ junctions each contributing $\sim 3.7$ trits
of ternary diversity. Multiplying by primary combinatorial diversity
($7{,}500 = 3^{k'}$ where $k' \approx 8$):
$N_{\mathrm{primary}} = 7{,}500 \times 2.6 \times 10^5 = 1.95 \times 10^9 \approx 2 \times 10^9$. \qed
\end{proof}

\begin{proposition}[Somatic Hypermutation as Richness Amplification]
\label{prop:shm}
Somatic hypermutation (SHM) in germinal centres amplifies antibody diversity
from $\sim 2 \times 10^9$ to $\sim 10^{11}$ variants, a factor of $\sim 50$,
by introducing point mutations at a rate of $\sim 10^{-3}$ per base per
division in the variable region. In partition terms, SHM increases
$\delta(\text{antibody})$ by exploring the sequence neighbourhood of each
primary variant, raising $\Rich(\text{antibody})$ from its primary value
to a higher, antigen-optimised value.
\end{proposition}

\begin{proof}
The variable region of an antibody is $\sim 360$ bp. At mutation rate
$\mu = 10^{-3}$ per base per division over $N = 5$ rounds of division in
the germinal centre, the expected number of mutations is
$360 \times 10^{-3} \times 5 = 1.8$ per clone per round. The number of
distinct mutated sequences reachable from a starting sequence in $N$ rounds
is approximately:
\begin{equation}
  N_{\mathrm{SHM}} \approx \binom{360}{2} \times 3^2 \approx 64{,}440 \times 9 \approx 580{,}000 \sim 10^{5.8}.
\end{equation}
Multiplied by the primary diversity ($2 \times 10^9 \approx 10^{9.3}$),
the total post-SHM diversity is $\sim 10^{9.3 + 5.8 - \text{(degeneracy correction)}}
\approx 10^{11}$, consistent with the empirical estimate.

The partition depth of the antibody repertoire increases by
$\Delta\Depth_{\mathrm{SHM}} = \log_2(10^{5.8}) \approx 19$ bits per primary
variant, expanding the richness $\Rich(\text{antibody})$ and enabling fine
discrimination within richness classes. \qed
\end{proof}

%------------------------------------------------------------------------------
\section{The Immune Cascade as Categorical Filter Sequence}
\label{sec:immune_cascade}
%------------------------------------------------------------------------------

\begin{definition}[Categorical Filter]
\label{def:cat_filter}
A \emph{categorical filter} is a continuous map $\Phi: \mathcal{C}_1 \to \mathcal{C}_2$
between categorical spaces that reduces the equivalence classes of its input
by a factor $\rho \geq 1$:
\begin{equation}
  |\mathcal{C}_2| = |\mathcal{C}_1| / \rho,
\end{equation}
where $\rho$ is the \emph{reduction ratio}. A categorical filter preserves
the categorical structure (it is a functor: morphisms in $\mathcal{C}_1$ map
to morphisms in $\mathcal{C}_2$) while collapsing $\rho$ input classes
into each output class.
\end{definition}

\begin{theorem}[Immune System as Three-Stage Filter Cascade]
\label{thm:immune_cascade}
The immune response is a three-stage categorical filter cascade, each stage
reducing categorical richness by a factor of $\sim 10^3$:
\begin{align}
  \text{Stage 1 (Innate):}   &\quad \Rich_{\mathrm{pathogen}} \sim 10^9 \xrightarrow{\Phi_{\mathrm{innate}}}
                              \Rich_{\mathrm{pattern}} \sim 10^6 \quad (\rho_1 \sim 10^3) \notag \\
  \text{Stage 2 (Adaptive):} &\quad \Rich_{\mathrm{pattern}} \sim 10^6 \xrightarrow{\Phi_{\mathrm{adaptive}}}
                              \Rich_{\mathrm{antigen}} \sim 10^3 \quad (\rho_2 \sim 10^3) \notag \\
  \text{Stage 3 (Memory):}   &\quad \Rich_{\mathrm{antigen}} \sim 10^3 \xrightarrow{\Phi_{\mathrm{memory}}}
                              \Rich_{\mathrm{eliminated}} \sim 10^0 \quad (\rho_3 \sim 10^3). \notag
\end{align}
Three orders of magnitude per stage, three stages, nine orders of magnitude
total: from complex pathogen ($\Rich \sim 10^9$) to eliminated (single
categorical state, $\Rich = 1$).
\end{theorem}

\begin{proof}
\textbf{Stage 1: Innate immunity.}
Pattern-recognition receptors (PRRs: TLRs, NLRs, RLRs) recognise
pathogen-associated molecular patterns (PAMPs). PAMPs are molecular patterns
that are conserved across broad classes of pathogens but absent (or richness-suppressed)
in self-proteins. TLR4 recognises LPS from Gram-negative bacteria (richness class:
$\Rich_{\mathrm{LPS}} \sim 10^6$). The PRR repertoire collectively covers
$\Rich$ classes from $10^6$ to $10^9$. Innate immunity is thus a filter from
the full pathogen richness range to the PAMP richness range:
$\rho_1 = 10^9/10^6 = 10^3$.

\textbf{Stage 2: Adaptive immunity.}
After innate recognition, dendritic cells (DCs) process and present antigens
via MHC. T-cell activation requires co-stimulation (CD28:CD80/86), restricting
the responding cells to those both antigen-specific and co-stimulated. This
narrows from PAMP richness class ($10^6$) to specific antigen richness class
($10^3$): $\rho_2 = 10^3$. B-cell somatic hypermutation and affinity maturation
further refine the response within the antigen richness class.

\textbf{Stage 3: Memory.}
Long-lived memory B and T cells encode the specific antigen coordinate as
a categorical fixed point in $\Sspace$: upon re-exposure, the memory response
is activation at $\Rich \sim 1$ (the pathogen is ``known'' and its categorical
richness in immune space is 1). $\rho_3 = 10^3/10^0 = 10^3$.

The product $\rho_1 \rho_2 \rho_3 = 10^9$ spans the full range from complex
unknown pathogen to recognised and eliminated threat. \qed
\end{proof}

\begin{remark}[Cross-Reactivity as Categorical Overlap]
Immune cross-reactivity (one antibody reacting with multiple antigens) is
categorical overlap: the two antigens share a richness class, so the filter
$\Phi_{\mathrm{adaptive}}$ maps them to the same output class. This is not a
failure of specificity; it is a feature of categorical filtering. Richness
class overlap is more common near the self-threshold, explaining why
cross-reactive autoimmunity arises preferentially with antigens near
$\Rich_{\mathrm{thresh}}$.
\end{remark}

%------------------------------------------------------------------------------
\section{MHC Binding Prediction and Therapeutic Implications}
\label{sec:mhc_prediction}
%------------------------------------------------------------------------------

\subsection{Predicting Immunogenicity from Categorical Richness}

The correlation $r = 0.73$ between $\Rich(\text{source})$ and MHC binding
affinity (Theorem~\ref{thm:mhc_filter}) constitutes a quantitative predictor
of immunogenicity.

\begin{proposition}[Richness-Based Immunogenicity Prediction]
\label{prop:immunogenicity}
For a target protein $T$ with calculated categorical richness $\Rich(T)$:
\begin{itemize}
  \item $\Rich(T) > \Rich_{\mathrm{thresh}} \approx 10^{4.5}$: protein has
        high MHC presentation probability (the source richness is in the range
        matched to MHC filter depth). Predicted to be an effective immunotherapy
        target if tumour-associated.
  \item $\Rich(T) < \Rich_{\mathrm{thresh}}$: protein has low MHC presentation
        probability. Not an effective immunotherapy target; immune evasion likely.
\end{itemize}
\end{proposition}

\begin{proof}
From Theorem~\ref{thm:mhc_filter}, MHC binding affinity is maximised for
peptides from proteins with $\Rich \approx \Rich_{\mathrm{self}} \approx 10^{4.5}$.
Proteins with $\Rich > 10^{4.5}$ are in the same richness range and present
peptides efficiently. Proteins with $\Rich < 10^{4.5}$ are outside the
MHC filter window and are presented poorly: their peptides have lower partition
depth than the groove's tuned $\Depth_{\mathrm{MHC}}$, reducing affinity.

The practical implication for immunotherapy: tumour antigens should be
selected from the high-$\Rich$ range. This is consistent with the empirical
observation that the most immunogenic neoantigens are derived from highly
connected (hub) proteins in the protein interaction network, which are
precisely the high-$N_{\downarrow}$ (high-$\Rich$) proteins. \qed
\end{proof}

\subsection{Vaccine Design Implications}

\begin{corollary}[Categorical Richness Criterion for Vaccine Antigens]
\label{cor:vaccine_design}
Effective vaccine antigens should have $\Rich_{\mathrm{antigen}} \in [10^{4.0}, 10^{5.0}]$:
sufficiently high to be presented by MHC (above $\Rich_{\mathrm{thresh}}$)
but sufficiently distinct from the mean self-richness to avoid tolerance
at the population level. Antigens with $\Rich \gg 10^5$ risk molecular mimicry
with highly connected host proteins.
\end{corollary}

\begin{proof}
The MHC binding window is approximately $[\Rich_{\mathrm{self}} - 2\sigma,
\Rich_{\mathrm{self}} + 2\sigma]$, where $\sigma \sim 0.5$ log units.
For efficient presentation, the antigen richness must fall within this window.
The lower bound $\Rich > 10^{4.0}$ ensures presentation above the 2$\sigma$
lower boundary. The upper bound $\Rich < 10^{5.0}$ prevents entry into the
richness range of highly connected host hub proteins, which could cross-react
via Corollary~\ref{cor:autoimmunity}. \qed
\end{proof}

%------------------------------------------------------------------------------
\section{Complement and the Categorical Opsonisation Cascade}
\label{sec:complement}
%------------------------------------------------------------------------------

The complement system provides an independent illustration of categorical
filtering. The C3 complement protein (the central component of all three
complement pathways) undergoes spontaneous hydrolysis at a basal rate,
generating C3b, which covalently attaches to nearby surfaces.

\begin{proposition}[Complement as Richness-Threshold Opsonisation]
\label{prop:complement}
Complement deposition is a categorical richness detector: C3b binds
indiscriminately to nearby surfaces, but the downstream effector response
(phagocytosis, MAC pore formation) is gated by the categorical richness of
the C3b-coated target. Self-surfaces express complement regulatory proteins
(CD55, CD59, factor~H) that reduce their effective richness below the
complement-activation threshold. Non-self surfaces lack these regulators,
maintaining high effective richness, and are opsonised.
\end{proposition}

\begin{proof}
Factor H binding (the key regulatory step) is mediated by polyanion
patterns on the surface. Self-surfaces display sialic acid and heparan
sulfate (high-charge-density polyanions) that recruit factor~H.
Factor~H binding reduces the effective C3b deposition rate below the
threshold for amplification by the alternative pathway.
In categorical terms: self-surfaces have their partition depth of
complement-activating patterns reduced by regulatory proteins to below
$\Depth_{\mathrm{comp}}$ (the complement-activation threshold). Non-self
surfaces lack this depth reduction and reach the amplification threshold. \qed
\end{proof}

%------------------------------------------------------------------------------
\section{Natural Killer Cells and Richness-Deficiency Detection}
\label{sec:nk_cells}
%------------------------------------------------------------------------------

Natural killer (NK) cells provide a complementary filter to T cells: they kill
cells that have \emph{lost} categorical richness markers (MHC class~I).

\begin{proposition}[NK Cell ``Missing Self'' as Richness-Deficiency Detection]
\label{prop:nk_missing_self}
NK cells kill targets that are below $\Rich_{\mathrm{thresh}}$ in the
``MHC-presentation richness'' dimension. Normal cells present MHC class~I
(self richness $> \Rich_{\mathrm{thresh}}$): NK cells are inhibited.
Virally infected or transformed cells downregulate MHC class~I to evade T cells,
reducing their MHC-presentation richness below $\Rich_{\mathrm{thresh}}$:
NK cells are activated.
\end{proposition}

\begin{proof}
NK cell inhibitory receptors (KIRs, NKG2A/CD94) bind MHC class~I molecules.
When MHC class~I is present at normal surface density (healthy cells),
inhibitory signalling dominates. When MHC class~I is downregulated
(virally infected, tumour cells), the inhibitory signal falls below threshold,
and activating receptors (NKG2D, NCR receptors recognising stress ligands)
drive cytotoxicity.

In categorical terms: inhibitory KIR signal is proportional to
$\Rich_{\mathrm{MHC}}$ on the target surface. Below $\Rich_{\mathrm{thresh}}$,
the NK cell classifies the target as richness-deficient (non-self) and kills.
This is the dual complement to T-cell activity: T cells kill targets that
exceed the self-richness window (foreign antigens); NK cells kill targets
that fall below the self-richness floor (MHC loss). Together, they enforce
an immune richness corridor around $\Rich_{\mathrm{self}} \approx 10^{4.5}$. \qed
\end{proof}

%==============================================================================
\section*{Chapter Summary}
\label{sec:ch23_summary}
%==============================================================================

\begin{chapterbox}[title=Chapter 23 --- Key Results Established]
\begin{enumerate}[label=\textbf{\arabic*.}]
  \item \textbf{Categorical Richness (Definition~\ref{def:categorical_richness})}:
        $\Rich(P) = \log_2 \delta(P) + \log_2 N_{\downarrow}(P)$.
        Partition depth plus downstream connectivity. Units: bits.

  \item \textbf{MHC as Ambiguous Filter (Theorem~\ref{thm:mhc_filter})}:
        MHC binding correlates with source protein $\Rich$ at $r = 0.73$,
        $p < 10^{-12}$. MHC is a depth-matched richness filter, not a sequence-
        specific lock.

  \item \textbf{Self/Non-Self Threshold (Theorem~\ref{thm:self_nonself})}:
        $\Rich_{\mathrm{thresh}} \approx 10^{4.5}$, derived from
        $N_{\mathrm{TCR}} \times \rho = 10^7 \times 10^{-2.5} = 10^{4.5}$.
        Above: tolerised (self). Below: attacked (non-self).

  \item \textbf{VDJ as $3^k$ Cascade (Theorem~\ref{thm:vdj_ternary})}:
        Each junctional position is a ternary choice (add/neutral/delete).
        Primary diversity $\approx 2 \times 10^9$; post-SHM diversity
        $\approx 10^{11}$.

  \item \textbf{Three-Stage Filter Cascade (Theorem~\ref{thm:immune_cascade})}:
        Innate $\to$ adaptive $\to$ memory, each reducing $\Rich$ by
        $\sim 10^3$. Total reduction: $10^9$ (pathogen) $\to$ $10^0$ (eliminated).

  \item \textbf{Richness-Based Immunogenicity (Proposition~\ref{prop:immunogenicity})}:
        Proteins with $\Rich > 10^{4.5}$ are efficiently MHC-presented and are
        effective immunotherapy targets.

  \item \textbf{Vaccine Design Criterion (Corollary~\ref{cor:vaccine_design})}:
        Optimal vaccine antigens have $\Rich \in [10^{4.0}, 10^{5.0}]$.

  \item \textbf{NK ``Missing Self'' (Proposition~\ref{prop:nk_missing_self})}:
        NK cells enforce the lower richness boundary; T cells enforce the upper.
        Together they maintain a self-richness corridor
        centred on $\Rich_{\mathrm{self}} \approx 10^{4.5}$.
\end{enumerate}
\end{chapterbox}
